\section{0x32\ 数学\ 集合幂级数}
\label{sec:0x32}

\subsection{FWT}

\Highlight{位运算卷积}：已知序列 $a_0, \dots, a_{2^m - 1}$ 和序列 $b_0, \dots, b_{2^m - 1}$，求出 $c_0, \dots, c_{2^m - 1}$，其中

\[
c_k = \sum_{i \oplus j = k} a_ib_j
\]
其中 $\oplus$ 表示一种位运算（or / and / xor）。

\subsubsection{DWT}

\begin{quoteblock}
FFT 多项式乘法的核心思想：$A, B$ 的系数转 $A, B$ 的点值，$A, B$ 的点值点积转 $C$ 的点值，$C$ 的点值转 $C$ 的系数。
\end{quoteblock}

我们想要 DWT 与 IDWT 有着类似的功能。设 $\mathrm{dwt}(a)$ 表示幂级数 $a$ 经过 DWT 变换后得到的幂级数（点值），我们希望

\[
a \times b = c \iff \mathrm{dwt}(a) \cdot \mathrm{dwt}(b) = \mathrm{dwt}(c)
\]
由于 DFT 是线性变换，我们希望 DWT 也是线性变换。

记序列长度 $n = 2^m$，设 $\mathrm{dwt}(a)_i = \sum_{j = 0}^{n - 1} c(i, j)a_j$，其中变换系数 $c(i, j)$ 表示变换前下标 $j$ 对变换后下标 $i$ 的贡献。

由 $\mathrm{dwt}(a) \cdot \mathrm{dwt}(b) = \mathrm{dwt}(c)$，得

\[
\begin{gathered}
\mathrm{dwt}(c)_i = \mathrm{dwt}(a)_i \cdot \mathrm{dwt}(b)_i \\
\left( \sum_{p = 0}^{n - 1}c(i, p)c_p \right) = \left( \sum_{j = 0}^{n - 1}c(i, j)a_j \right) \cdot \left( \sum_{k = 0}^{n - 1}c(i, k)b_k \right) \\
\sum_{p = 0}^{n - 1}c(i, p)c_p = \sum_{j = 0}^{n - 1}\sum_{k = 0}^{n - 1}c(i, j)c(i, k)a_jb_k
\end{gathered}
\]
由 $a \times b = c$，得

\[
\begin{aligned}
\sum_{p = 0}^{n - 1} c(i, p)c_p & = \sum_{p = 0}^{n - 1} c(i, p) \sum_{t_1 \oplus t_2 = p} a_{t_1}b_{t_2} \\
& = \sum_{p = 0}^{n - 1} \sum_{t_1 \oplus t_2 = p} c(i, t_1 \oplus t_2) a_{t_1}b_{t_2} \\
& = \sum_{t_1 = 0}^{n - 1} \sum_{t_2 = 0}^{n - 1} c(i, t_1 \oplus t_2) a_{t_1}b_{t_2}
\end{aligned}
\]
对比上下两段的最后一个式子，\textbf{发现充要条件为} $c(i, j)c(i, k) = c(i, j \oplus k)$。

由于位运算可以\textbf{按位考虑}，可以这样构造 $c(i, j)$：\textbf{先得到单个二进制位的转移系数} $c([0, 1], [0, 1])$，满足 $c(i, j)c(i, k) = c(i, j \oplus k)$。对于所有的 $c(i, j)$，我们将 $i, j$ 二进制分解，\textbf{将对应位的转移系数相乘}

\[
c(i, j) = \prod_{x = 0}^{m - 1} c(i_x, j_x)
\]
（其中 $i_x, j_x$ 分别表示 $i, j$ 的二进制下第 $x$ 位）

代入条件 $c(i, j)c(i, k) = c(i, j \oplus k)$，可验证其成立

\[
\begin{aligned}
c(i, j)c(i, k) & = \left(\prod_{x = 0}^{m - 1}c(i_x, j_x)\right) \left( \prod_{x = 0}^{m - 1}c(i_x, k_x) \right) \\
& = \prod_{x = 0}^{m - 1} c(i_x, j_x)c(i_x, k_x) \\
& = \prod_{x = 0}^{m - 1} c(i_x, j_x \oplus k_x) \\
& = c(i, j \oplus k)
\end{aligned}
\]
现在考虑优化计算 DWT。考虑\textbf{按最高位拆分}，记 $i', j'$ 表示 $i, j$ 去掉最高位的值，记 $i_h$ 表示 $i$ 的最高位

\[
\begin{aligned}
\mathrm{dwt}(a)_i & = \sum_{j = 0}^{n - 1}c(i, j)a_j \\
& = \sum_{j = 0}^{n / 2 - 1}c(i, j)a_j + \sum_{j = n / 2}^{n - 1}c(i, j)a_j \\
& = c(i_h, 0) \left( \sum_{j = 0}^{n / 2 - 1}c(i', j')a_j \right) + c(i_h, 1)\left( \sum_{j = n / 2}^{n - 1}c(i', j')a_j \right)
\end{aligned}
\]
发现括号里的式子是一个子问题。记 $a_0$ 表示幂级数 $a$ 最高位为 $0$ 的部分，$a_1$ 表示幂级数 $a$ 最高位为 $1$ 的部分。

\begin{itemize}
\item 对于 $k < n / 2$，代入 $i = k$ 得

\[
\mathrm{dwt}(a)_k = c(0, 0) \cdot \mathrm{dwt}(a_0)_k + c(0, 1) \cdot \mathrm{dwt}(a_1)_k
\]

\item 对于 $k < n / 2$，代入 $i = k + n / 2$ 得

\[
\mathrm{dwt}(a)_{k + n / 2} = c(1, 0) \cdot \mathrm{dwt}(a_0)_k + c(1, 1) \cdot \mathrm{dwt}(a_1)_k
\]
\end{itemize}

如果我们知道 $\mathrm{dwt}(a_0)_i, \mathrm{dwt}(a_1)_i$ 在 $i = 0, 1, \dots, n / 2 - 1$ 处的值，就可以 $\mathcal{O}(n)$ 求出 $\mathrm{dwt}(a)_i$ 在 $i = 0, 1, \dots, n - 1$ 处的值。

这是一个分治的过程，不断分治直到仅剩一个项即可。

时间复杂度 $\mathcal{O}(m2^m)$。

\subsubsection{IDWT}

在 DWT 中，我们只需要

\[
c = \begin{bmatrix}
c(0, 0) & c(0, 1) \\
c(1, 0) & c(1, 1)
\end{bmatrix}
\]
这四个数即可完成变换，记该矩阵为位矩阵 $c$。

如果我们要进行 IDWT，则需要位矩阵 $c$ 的逆矩阵 $c^{-1}$。类似地，有

\[
a_i = \sum_{j = 0}^{n - 1} c^{-1}(i, j)\mathrm{dwt}(a)_j
\]
套用 DWT 的模板即可。

在构造矩阵 $c$ 的时候，我们需要小心一点，\textbf{得确保矩阵 $c$ 存在逆矩阵 $c^{-1}$}。

\subsubsection{位运算卷积}

\Highlight{or 卷积}：or 运算的位矩阵，及其逆矩阵

\[
c = \begin{bmatrix}
1 & 0 \\
1 & 1
\end{bmatrix}, \quad c^{-1} = \begin{bmatrix}
1 & 0 \\
-1 & 1
\end{bmatrix}
\]
or 卷积的正变换等价于\textbf{子集求和}，因为 $c(i, j) = [i \operatorname{and} j = j]$。

\Highlight{and 卷积}：and 运算的位矩阵，及其逆矩阵

\[
c = \begin{bmatrix}
1 & 1 \\
0 & 1
\end{bmatrix}, \quad c^{-1} = \begin{bmatrix}
1 & -1 \\
0 & 1
\end{bmatrix}
\]
and 卷积的正变换等价于\textbf{超集求和}，因为 $c(i, j) = [i \operatorname{or} j = j]$。

\Highlight{xor 卷积}：xor 运算的位矩阵，及其逆矩阵

\[
c = \begin{bmatrix}
1 & 1 \\
1 & -1
\end{bmatrix}, \quad c^{-1} = \frac{1}{2}\begin{bmatrix}
1 & 1 \\
1 & -1
\end{bmatrix}
\]
xor 卷积的正变换，其变换系数 $c(i, j) = (-1)^{\mathrm{popcount}(i \operatorname{and} j)}$。

xor 卷积的 IDWT 相较于 DWT，每一位的变换系数需要乘以 $\frac{1}{2}$，一共有 $m$ 位，总共的影响为 $\frac{1}{2^m} = \frac{1}{n}$。所以 xor 卷积的 IDWT 只需套用 DWT 的模板，最后乘以 $\frac{1}{n}$ 即可。

\InputCode{cpp}{contents/0x30 数学/code/0x32 FWT.cpp}

\subsection{子集 dp}

\Highlight{子集 dp}：已知序列 $a_0, \dots, a_{2^m - 1}$，求出 $b_0, \dots, b_{2^m - 1}$，其中 $b_i = \sum_{j \subseteq i} a_j$。

将一个二进制状态视作一个 $m$ 维的坐标，第 $i$ 维的取值等于二进制下第 $i$ 位的值。则二进制下的子集和，相当于是要求\textbf{每个维度都小于等于当前坐标}的所有坐标的权值和，运用\textbf{高维前缀和}即可。

当信息满足可减性时，可以进行子集 dp 的逆过程。

时间复杂度 $\mathcal{O}(m2^m)$。

\begin{minted}{cpp}
// 注意：要先枚举维度 i，再枚举状态 S
for (int i = 0; i < m; i ++) {
    for (int S = 0; S < (1 << m); S ++) {
        if (S >> i & 1) {
            f[S] += f[S ^ (1 << i)];
        }
    }
}
\end{minted}

\Highlight{半在线子集 dp}：求出 $b_0, \dots, b_{2^m - 1}$，其中 $b_i$ 与 $\sum_{j \subsetneqq i} b_j$ 强相关。

按照 $|S|$ 从小到大分层计算。每一层都做一次子集 dp，然后更新当前层的所有状态。

时间复杂度 $\mathcal{O}(m^22^m)$。

\begin{minted}{cpp}
for (int s = 1; s <= m; s ++) {
    for (int S = 0; S < (1 << m); S ++) {
        tmp[S] = f[S];
    }
    for (int i = 0; i < m; i ++) {
        for (int S = 0; S < (1 << m); S ++) {
            if (S >> i & 1) {
                tmp[S] += tmp[S ^ (1 << i)];
            }
        }
    }
    
    for (int S = 0; S < (1 << m); S ++) {
        if (__builtin_popcount(S) == s) {
            // 使用 tmp[S] 更新 f[S]
        }
    }
}
\end{minted}

\subsection{子集卷积}

\Highlight{子集卷积}：已知序列 $a_0, \dots, a_{2^m - 1}$ 和序列 $b_0, \dots, b_{2^m - 1}$，求出 $c_0, \dots, c_{2^m - 1}$，其中

\[
c_k = \sum_{i \operatorname{or} j = k, i \operatorname{and} j = \varnothing} a_ib_j
\]
注意到 $i \operatorname{and} j = \varnothing$ 的充要条件为 $|i| + |j| = |i \operatorname{or} j|$。于是引入序列 $f$ 的占位多项式 $F$

\[
F_{p, i} = \begin{cases}
f_i & |i| = p\\
0 & |i| \neq p
\end{cases}
\]
将占位多项式代入到子集卷积的式子中

\[
C_{s, k} = \sum_{p + q = s} \sum_{i \operatorname{or} j = k}A_{p, i}B_{q, j}
\]
从 or 卷积的角度来看，有

\[
C_s = \sum_{p + q = s}A_p \operatorname{or} B_q
\]
于是将 $A_p$ 和 $B_q$ 进行 DWT，求出所有的 $C_s$ 后再进行 IDWT 即可。

时间复杂度 $\mathcal{O}(m^22^m)$，空间复杂度 $\mathcal{O}(m2^m)$。

\InputCode{cpp}{contents/0x30 数学/code/0x32 子集卷积.cpp}

\Highlight{半在线子集卷积}：已知序列 $g_0, \dots, g_{2^m - 1}$，求出 $f_0, \dots, f_{2^m - 1}$，其中对于 $k > 0$ 满足

\[
f_k = \sum_{i \operatorname{or} j = k, i \operatorname{and} j = \varnothing} f_ig_j
\]
先根据具体问题的定义得出 $f_0$（$g_0$ 一般来说等于 $0$）。

然后按照 $|S|$ 从小到大分层计算，每一层都根据前面层的 DWT 结果得到当前层的 DWT 结果，然后先 IDWT 回去，更新当前层的所有状态，然后再 DWT 回来，为后面层的计算做准备。

时间复杂度 $\mathcal{O}(m^22^m)$，空间复杂度 $\mathcal{O}(m2^m)$。

\begin{minted}{cpp}
// 半在线子集卷积
std::vector<int> subset_conv(std::vector<int> g) {
    int n = g.size(), m = std::__lg(n);
    std::vector<std::vector<int>> F(m + 1, std::vector<int>(n, 0));
    std::vector<std::vector<int>> G(m + 1, std::vector<int>(n, 0));

    for (int i = 1; i < n; i ++) { // 一般来说 g_0 = 0，所以不考虑 G[0]
        G[__builtin_popcount(i)][i] = g[i];
    }
    for (int i = 1; i <= m; i ++) {
        dwtOr(G[i], +1);
    }

    int f0 = 1;
    std::fill(F[0].begin(), F[0].end(), f0); // 一般来说 f_0 = 1，这一步模拟了 dwtOr(F[0], +1)

    std::vector<int> f(n);
    f[0] = f0;
    for (int i = 1; i <= m; i ++) {
        for (int j = 1; j <= i; j ++) {
            for (int k = 0; k < n; k ++) {
                add(F[i][k], 1ll * G[j][k] * F[i - j][k] % mod);
            }
        }
        dwtOr(F[i], -1);

        for (int S = 0; S < n; S ++) {
            if (__builtin_popcount(S) == i) {
                // 在此更新 F[i][S]，即为真正的 f[S]
                f[S] = F[i][S];
            } else {
                F[i][S] = 0;
            }
        }
        dwtOr(F[i], +1);
    }
    return f;
}
\end{minted}

